\documentclass[11pt]{article}
\usepackage{fontspec}
\usepackage[french]{babel}
\usepackage{amsmath}
\usepackage{amssymb}
\usepackage{amsthm}
\usepackage{mathrsfs}
\usepackage{tikz-cd}
\usepackage{enumitem}

\newcommand{\printedpage}[1]{\marginpar{\footnotesize p.~#1}}
\newcommand{\msp}[1]{\marginpar{\footnotesize [#1]}}
\newcommand{\Hom}{\underline{\mathrm{Hom}}}
\newcommand{\Fib}{\underline{\mathrm{Fib}}}
\newcommand{\Cat}{\mathrm{Cat}}
\newcommand{\HOT}{\mathrm{HOT}}
\newcommand{\Hot}{\mathrm{Hot}}
\newcommand{\hot}{\mathrm{hot}}
\newcommand{\id}{\mathrm{id}}

\begin{document}

\printedpage{130}
% head: A. Grothendieck : Dérivateurs, Ch. XII, Édition des Universités de Montpellier II et Paris VII

\[
\begin{tikzcd}[column sep=small, row sep=small]
 & C(X') & \\
C(X) \arrow[ur, "C(i')"] & & X' \arrow[ul, "i_{X'}"'] \\
 & X \arrow[ul, "i_X"] \arrow[ur, "i'"'] &
\end{tikzcd}
\]

est commutatif (par fonctorialité de $i_X : X \longrightarrow C(X)$), et on applique ceci au cas de $X' = C'(X)$, $i' = i'_X$.

Cette remarque nous encourage à faire un saut dans l'abstraction, et d'introduire la catégorie\msp{page 198}
\[
\mathcal{C} = \mathcal{C}_S = \begin{cases}
\text{catégorie des couples } (C,i), \text{ avec} \\
C : \Cat/S \longrightarrow \Cat/S \text{ se factorisant par } \Fib_W S, \\
i : \id_{\Cat/S} \longrightarrow C \text{ tel que } i_X \in W \ \forall\ X \text{ dans } \Cat/S,
\end{cases}
\tag*{$(*)$}
\]
une flèche
\[
(C,i) \longrightarrow (C',i')
\]
étant une flèche
\[
C \xrightarrow{u} C' \qquad \text{(dans } \Hom(\Cat/S, \Cat/S))
\]
rendant commutatif le carré \texttt{[lire : le diagramme]}
\[
\begin{tikzcd}[column sep=small]
 & \id_{\Cat/S} \arrow[dl, "i"'] \arrow[dr, "i'"] & \\
C \arrow[rr, "u"] & & C'.
\end{tikzcd}
\]

En d'autres termes, $\mathcal{C}$ est défini comme une sous-catégiorie pleine de % [sic] sous-catégiorie -> sous-catégorie, misspelled word
\[
\id_{\Cat/S}\backslash\Hom(\Cat/S, \Cat/S),
\]
à savoir celle dont les objets $(C,i)$ satisfont les deux conditions $(*)$ plus haut (impliquant l'une et l'autre le localiseur $W$).

\msp{page 199}Ces conditions impliquent ceci

\begin{enumerate}
\item[a)] $\forall\ X$ dans Cat/$S$, $(u_X : C(X) \longrightarrow C'(X)) \in W$, puisque
\[
\begin{tikzcd}[column sep=small]
 & X \arrow[dl, "i_X\in W"'] \arrow[dr, "i'_X\in W"] & \\
C(X) \arrow[rr, "u_X"] & & C'(X),
\end{tikzcd}
\]
\printedpage{131}
% head: A. Grothendieck : Dérivateurs, Ch. XII, Édition des Universités de Montpellier II et Paris VII
et

\item[b)] si $W$ satisfait W(4b), de sorte que $u_X \in W^{\mathrm{f}}_S$, alors ce définit un \textit{isomorphisme} entre $\varphi_C(X)$ et $\varphi_{C'}(X)$ % [sic] ce définit -> ceci définit, missing word
\[
\begin{tikzcd}[row sep=small]
s \arrow[r, mapsto] & \hot_W(C(X)_s) \arrow[d, "\wr\, u_{X,s}"] \\
s \arrow[r, mapsto] & \hot_W(C'(X)_s).
\end{tikzcd}
\]
\end{enumerate}

Mieux encore, dans ce cas, le foncteur composé
\[
\mathcal{C} \times \Cat/S \longrightarrow \Fib_W S \longrightarrow \HOT^{\mathrm{lc}}_W(S)
\]
donne naissance à
\[
\mathcal{C} \longrightarrow \Hom(\Cat/S, \HOT^{\mathrm{lc}}_W(S))
\]
qui transforme toute flèche de $\mathcal{C}$ en \texttt{[un]} \textit{isomorphisme,} donc se factorise par le groupoïde fondamental de $\mathcal{C}$,
\[
\Pi\mathcal{C} \longrightarrow \Hom(\Cat/S, \HOT^{\mathrm{lc}}_W(S)).
\]

\msp{page 200}Quant au foncteur envisagé tantôt
\[
\begin{aligned}
\mathcal{C} &\quad\longrightarrow\quad \underbrace{\Hom\,(\Cat/S, \Hom(S^\circ, \Hot_W))}_{\simeq\Hom((\Cat/S)\times S^\circ, \Hot_W)} \\
C &\quad\longmapsto\quad ((X,s) \longmapsto \hot_W(C(X)_s))\,,
\end{aligned}
\]
il se déduit du précédent en composant avec
\[
\begin{tikzcd}[column sep=small]
\HOT_W(S) \arrow[rr] \arrow[dr, dotted] & & \Hom(S^\circ, \Hot_W), \\
 & \Hom(\Pi S^\circ, \Hot_W) \arrow[ur, dotted] &
\end{tikzcd}
\]
plus exactement, avec le foncteur qui s'en déduit
\[
\Hom(\Cat/S, \HOT_W(S)) \longrightarrow \Hom(\Cat/S, \Hom(S^\circ, \Hot_W)).
\]

On voit poindre ici un principe plausible de canonicité d'isomorphismes, en ce qui concerne les foncteurs
\[
C_W : \Cat/S \longrightarrow \HOT^{\mathrm{lc}}_W(S)
\]
définis par les $C \in \mathrm{Ob}\,\mathcal{C}$, par la relation
\[
\Pi\mathcal{C} \approx \ \text{catégorie ponctuelle, i.e. } \mathcal{C} \text{ est 1-connexe.}
\]

De plus, les arguments utilisés jusqu'ici sont de nature à tel point formelle, que les aspects spécifiques du Modéliseur Cat
\printedpage{132}
% head: A. Grothendieck : Dérivateurs, Ch. XII, Édition des Universités de Montpellier II et Paris VII
\msp{page 201}semblent n'y jouer pratiquement aucun rôle, ou un rôle très effacé.

Le carré de la page 197 montre surtout que la catégorie $\mathcal{C}$ est filtrante - $(C,i)$, $(C,i')$ s'envoient l'un et l'autre dans $(C'',i'')$, avec $C'' = C \circ C'$, et $i'' = (C * i') \circ i = (i * C') \circ i'$. Il y a même une loi de composition dans $\mathcal{C}$, nullement commutative ni unitaire c'est sûr, mais sûrement associative, et (si cette composition est notée $*$) avec des morphismes fonctoriels % [sic] (C,i') -> (C',i'), the second couple of the following formula is (C',i')
\[
\xi \xrightarrow{\alpha_{\xi,\xi'}} \xi * \xi' \xleftarrow{\beta_{\xi,\xi'}} \xi'
\]
qui assurent le caractère filtrant croissant de $\mathcal{C}$. Ces flèches semblent remplacer en quelque sorte l'unité bilatère absente, et on s'attend à ce qu'elles aient les mêmes excellentes propriétés que si on avait une telle unité, concernant les composés doubles $(\xi * \xi') * \xi'' \simeq \xi * (\xi' * \xi'')$,

\msp{page 202}
\[
(\xi * \xi') * \xi'' \xrightarrow{\sim} \xi * (\xi' * \xi'').
\]

Il y a deux façons d'envoyer un produit binaire
\[
\xi * \xi', \qquad \xi' * \xi'', \qquad \xi * \xi''
\]
dans le produit ternaire, et bien elles doivent dans chacun des trois cas, être égales. Et de même il y a, ceci admis, encore deux façons d'envoyer $\xi$ dans le produit ternaire, en factorisant soit par $\xi * \xi'$, soit par $\xi * \xi''$, et de même pour les deux autres foncteurs $\xi'$, $\xi''$; et là encore, ces deux façons doivent être égales, dans chacun de ces trois cas - donc $3 + 3 = 6$ égalités en tout à écrire et à vérifier - on a vu pire!

La philosophie calculatoire devrait être celle-ci : Si on a un produit $n$-uple
\[
\xi_1 * \xi_2 * \cdots * \xi_n
\]
(où je me dispense de mettre des parenthèses \dots), pour tout produit $m$-uple partiel $(m \in [1, n-1])$\msp{page 203}
\[
\xi_{i_1} * \xi_{i_2} * \cdots * \xi_{i_m} \qquad\qquad (i_1 < i_2 < \cdots < i_m),
\]
en utilisant à tour de bras les isomorphismes d'associativité et les $\alpha_{\xi,\xi'}$, $\beta_{\xi,\xi'}$, il n'y a pourtant qu'une seule façon de l'envoyer dans le produit $n$-uple. J'ai bien l'impression que ça doit être vrai. (Est-ce que ça implique déjà que $\mathcal{C}$ est 1-connexe?)

En fait, il y a bien une unité bilatère, mais elle n'est pas dans $\mathcal{C}$, mais dans la catégorie plus grande
\[
\id_{\Cat/S}\backslash\Hom(\Cat/S, \Cat/S),
\]
qui elle, comme toute catégorie
\[
\id_{\mathcal{X}}\backslash\Hom(\mathcal{X}, \mathcal{X}),
\]
\printedpage{133}
% head: A. Grothendieck : Dérivateurs, Ch. XII, Édition des Universités de Montpellier II et Paris VII
est une catégorie monoïdale, i.e. avec produit associatif \textit{et} unitaire. Les flèches $\alpha_{\xi,\xi'}$, $\beta_{\xi,\xi'}$ sont les flèches qui résultent de la présence de cette unité (vu que $\mathcal{C}$ est une sous-catégorie pleine), et alors il doit devenir clair qu'on a les compatibilités dites.

\msp{page 204}Mais même une catégorie monoïdale, est elle 1-connexe? Par exemple $B_G$, pour $G$ un groupe disons, est bien une catégorie monoïdale et même ‘groupale’, une $G$-catégorie (tout ce qu'il y a de strict), et pourtant elle n'est \textit{pas} 1-connexe!

Bien sûr, on peut court-circuiter la question de la 1-connexité de $\mathcal{C}$, en notant qu'après tout on a déjà a priori un objet privilégié dans
\[
\Hom(\Cat/S, \HOT^{\mathrm{lc}}_W(S)),
\]
i.e. un foncteur canonique
\[
\Cat/S \longrightarrow \HOT^{\mathrm{lc}}_W(S),
\]
savoir le foncteur composé
\[
\Cat/S \underbrace{\longrightarrow}_{\substack{\text{localisation}\\ \text{par } W_S}} (\Cat/S)W_S^{-1} \longrightarrow \underbrace{\HOT^{\mathrm{lc}}_W(S)}_{\overset{\text{déf}}{=} (\Fib_W S)(W^{\mathrm{f}}_S)^{-1}} \quad ,
\tag*{$(*)$}
\]
où la deuxième flèche est définie comme le quasi-inverse de
\[
(\Fib_W S)(W^{\mathrm{f}}_S)^{-1} \longrightarrow (\Cat/S)W_S^{-1}
\tag*{$(**)$}
\]

\msp{page 205}déduit de l'inclusion
\[
(\Fib_W S, W^{\mathrm{f}}_S) \hookrightarrow (\Cat/S, W_S).
\]

Le fait que $(*)$ soit une équivalence, se prouve justement par le fait de l'existence d'un foncteur $C : \Cat/S \longrightarrow \Cat/S$ et d'un $i : \id_{\Cat/S} \longrightarrow C$, ayant les propriétés qu'on sait, i.e. du fait que $C \neq \emptyset$. De plus, la donnée d'un objet $C$ de $\mathcal{C}$ donne un moyen commode d'expliciter le quasi-inverse dudit foncteur. Mais ce quasi-inverse a une existence canonique et indépendante du choix d'un $C$, i.e. d'un mode de calcul. Le fait que le foncteur composé $\varphi$ dans $(*)$ puisse se ‘relever’ (à isomorphisme près) en un foncteur % [sic] C \neq \emptyset -> \mathcal{C} \neq \emptyset, the non-emptiness concerns the category \mathcal{C}
\[
C : \Cat/S \longrightarrow \Fib_W S
\]
entre $\varphi(X) = \hot_{W,S}(X)$ et $\hot_{W,S}(C(X))$ et que l'isomorphisme puisse s'expliciter, lui aussi, par une flèche $X$-fonctorielle
\[
i_X : X \longrightarrow C(X), \qquad i_X \in W,
\]
dans Cat/$S$ lui-même (et non seulement dans le localisé $(\Cat/S)W_S^{-1}$),
\printedpage{134}
% head: A. Grothendieck : Dérivateurs, Ch. XII, Édition des Universités de Montpellier II et Paris VII
\msp{page 206}est une vertu particulière du \texttt{[mot illisible]}, qui ne doit pas pour autant faire oublier le caractère intrinsèque des flèches $(*)$, indépendamment de tout mode de calcul.

\marginpar{$\|$}% double vertical bar printed in the left margin beside the first lines of this paragraph
La clef de la situation, et des questions de canonicité, me semble être dans le fait que $(**)$ est une équivalence de catégories - et il faut avouer que, tout familier et quasiment banal qu'il soit finalement devenu pour moi, c'est là un fait tout à fait remarquable et nullement évident a priori. (Mais sans doute peut-on l'avoir aussi par les fourbis généraux à la \textsc{Quillen-Thomason}, sans se battre avec un formalisme de chemins \dots) Les ‘isomorphismes canoniques’ qu'on peut établir entre une ‘théorie $C$ ’ et une ‘théorie $C'$ ’ (via deux constructions différentes, p. ex., de catégories de chemins, ou via l'approche de \textsc{Quillen-Thomason}, sans chemins) doivent toujours revenir à ceux

\msp{page 207}qu'on peu définir via la comparaison de la ‘théorie $C$ ’ d'une part, la ‘théorie $C'$ ’ de l'autre, avec ce que nous donne directement la flèche composée $(*)$ (p. 204). Finalement, il me semble que \textit{l'existence} d'une théorie des ‘fibres homotopiques’ (ou ‘$W$-homotopiques’) d'une flèche $X \longrightarrow S$ dans Cat, résulte de l'équivalence $(**)$ (p. 204), et cette théorie pourrait sans doute dans une certaine mesure se développer à partir de là, sans choix d'un foncteur $C$ ni d'une structure de catégorie de modèles au sens de \textsc{Quillen}. % [sic] qu'on peu définir -> qu'on peut définir, verb form

\section*{4\quad Axiomatisation des ‘fibres $W$-homotopiques’.}

La reflexion de hier me fait revenir sur l'axiomatique dans des ‘catégories de modèles’ % [sic] reflexion -> réflexion, missing accent
\[
(\mathcal{M}, W)
\]
(\footnote[119]{$W$ satisfait l'axiome de saturation faible : les isomorphismes \texttt{[sont]} dans $W$, et si deux parmi $f$, $g$, $gf$ \texttt{[sont]} dans $W$, le troisième aussi.}) plus générales que Cat, dans l'esprit de VI (cf. théorème-scholie 2, p. 108-113). Pour le moment, on va supposer seulement $\mathcal{M}$ stable par produits fibrés (on n'a pas besoin d'objet final). On introduit l'ensemble
\[
\mathrm{Fib}_W \subseteq \mathrm{Fl}(\mathcal{M})
\]

\msp{page 208}des \textit{$W$-fibrations} (sous-entendu ‘ultrafortes’!) par la condition

\begin{quote}
$f : X \longrightarrow Y$ est dans $\mathrm{Fib}_W$ si et seulement si le foncteur changement de base par $X$,
\[
\begin{aligned}
Y' &\quad\longmapsto\quad X' = X \times_Y Y' \\
\mathcal{M}/Y &\quad\longrightarrow\quad \mathcal{M}/X
\end{aligned}
\]
\end{quote}

\end{document}
